\section*{Capítulo \ref{ch:b3:schrodinger-three-dimensions} --- La ecuación de Schrödinger en tres dimensiones}

\begin{solution}{exo:b3:schrodinger-three-dimensions:1}
(a) $\int|\psi|^2\dd^3r = |A|^2\,4\pi\int_0^\infty r^2\eu^{-r^2/2
\sigma^2}\dd r = |A|^2(2\pi\sigma^2)^{3/2}$: $A = (2\pi\sigma^2)^{
-3/4}$. (b) Por isotropía, $\langle r^2\rangle = 3\langle x^2\rangle =
3\sigma^2$, luego $\Delta x = \sigma$. (c) $\vect\jmath = \vect 0$: la
función de onda es real. (d) $\langle\vect p\rangle = \hbar\vect k_0$
y $\vect\jmath = \rho\,\hbar\vect k_0/m$: la misma nube, ahora a la
deriva.
\end{solution}

\begin{solution}{exo:b3:schrodinger-three-dimensions:2}
$3\,(1)$, $6\,(3)$, $9\,(3)$, $11\,(3)$, $12\,(1)$, $14\,(6)$ ---
y luego nada hasta $17$. Nivel $27$: $(3,3,3)$ y las permutaciones
de $(5,1,1)$ --- cuatro estados, ya que $27 = 9+9+9 = 25+1+1$. La simetría
cúbica solo permuta ejes, y ninguna permutación enlaza $(3,3,3)$ con
$(5,1,1)$: la coincidencia es aritmética (dos representaciones como suma de
tres cuadrados), no geométrica.
\end{solution}

\begin{solution}{exo:b3:schrodinger-three-dimensions:3}
(a) $\vect\jmath = |A|^2\hbar k/m\,\vect e_x$. (b) Nula: función
real. (c) $j = (|A|^2\hbar k/m)(1 - r^2)$: flujo incidente menos
flujo reflejado --- los términos cruzados se cancelan. (d) $j = \rho\,\hbar
m/Mb$ con $\rho = 1/2\pi$ (por unidad de ángulo): una circulación
permanente, el antepasado microscópico de las corrientes persistentes.
\end{solution}

\begin{solution}{exo:b3:schrodinger-three-dimensions:4}
El coeficiente es $\pi^2(\hbar c)^2/2\mu c^2 = \qty{3.76}{eV.nm^2}$.
(a) En $R = \qty{3}{nm}$: \qty{0.42}{eV}. (b) $E = \qty{2.16}{eV}$:
$\lambda = 1240/2.16 = \qty{575}{nm}$, verde amarillento. (c) Menor $R$,
mayor $E$: hacia el azul. (d) En el material masivo, el término de
confinamiento se anula: banda prohibida fija de \qty{1.74}{eV},
$\lambda = \qty{713}{nm}$, el mismo rojo profundo sea cual sea la muestra.
\end{solution}

\begin{solution}{exo:b3:schrodinger-three-dimensions:5}
(a) Univaluación: $\psi(\varphi + 2\pi) = \psi(\varphi)$ obliga a
$\eu^{2\pi\iu m} = 1$. (b) $E_m = \hbar^2m^2/2Mb^2$; los estados $\pm
m$ circulan en sentidos opuestos y se intercambian por reflexión especular
(o por inversión temporal) --- una simetría del anillo, de ahí la
degeneración. (c) La espira de corriente $I = q\hbar m/2\pi Mb^2$ lleva el
momento magnético $\mu = I\pi b^2 = m\,q\hbar/2M$ --- cuantizado en pasos de
$q\hbar/2M$, el patrón del magnetón de Bohr del
\cref{ch:b3:quantum-angular-momentum}.
(d) $\hbar^2/2m_{\text{e}}b^2 = \qty{1.9}{eV}$; los seis electrones
llenan $m = 0, \pm1$, y la primera excitación $\pm1 \to \pm2$ cuesta
$(4-1) \times 1.9 \approx \qty{5.8}{eV}$, es decir, $\lambda \approx
\qty{210}{nm}$ --- el vecindario ultravioleta correcto a partir de un anillo
y de nada más.
\end{solution}

\begin{solution}{exo:b3:schrodinger-three-dimensions:6}
(a) Elíjanse $n_x = 0..n$ y $n_y = 0..n - n_x$: $\sum_{n_x}(n - n_x
+ 1) = (n+1)(n+2)/2$. (b) Con espín: $2, 6, 12, 20$; acumulados $2,
8, 20, 40$. (c) $2$, $8$ y $20$ son mágicos; el fallo posterior dice que
el pozo nuclear no es armónico --- tiene el fondo más plano y, sobre todo,
un fuerte acoplamiento espín--órbita que rebaraja las capas hasta
$28, 50, 82$. (d) La estructura de capas exige niveles \emph{agrupados}
separados por huecos; los de la caja se reparten irregularmente sin huecos
grandes, de modo que no emerge ningún patrón de estabilidad nuclear.
\end{solution}

\begin{solution}{exo:b3:schrodinger-three-dimensions:7}
(a) $u_n = \sqrt{2/R}\sin(n\pi r/R)$, $E_n = n^2\pi^2\hbar^2/2MR^2$.
(b) $E_1 = \pi^2(\hbar c)^2/2Mc^2R^2 = 384210/(2 \times 939 \times
25) = \qty{8.2}{MeV}$. (c) La escala correcta: los nucleones en los núcleos
son estados cuánticos espaciados en MeV, como muestran sus espectros. (d)
$|u_1|^2$ alcanza su máximo en $r = R/2$: la probabilidad radial es máxima a
media distancia (el $r^2$ del elemento de volumen, escondido en
$u = r\varphi$, desplaza el máximo del centro), mientras que el
$|\varphi|^2$ de la caja unidimensional lo alcanza en el centro.
\end{solution}

\begin{solution}{exo:b3:schrodinger-three-dimensions:8}
(a) Hágase $N = 2N_{\text{estados}}(E_{\text{F}})$ e inviértase el recuento.
(b) Cobre: $E_{\text{F}} = \qty{7.1}{eV}$. (c) $v_{\text{F}} =
\sqrt{2E_{\text{F}}/m_{\text{e}}} = \qty{1.6e6}{m/s}$;
$T_{\text{F}} = E_{\text{F}}/k_{\text{B}} \approx \qty{8e4}{K}$. (d)
Los electrones son fermiones idénticos: el principio de exclusión de Pauli
(\cref{ch:b3:identical-particles}) admite como mucho dos por estado
orbital, de modo que la multitud debe apilarse hasta energías de
electronvoltios.
\end{solution}

\begin{solution}{exo:b3:schrodinger-three-dimensions:9}
(a) $k = \sqrt{2M(V_0 - |E|)}/\hbar$, $\kappa = \sqrt{2M|E|}/\hbar$;
la continuidad de $u'/u$ en $R$ da $k\cot kR = -\kappa$. (b) Cuando $|E|
\to 0$, $\kappa \to 0$: $\cot kR = 0$, $kR = \pi/2$, y entonces $V_0 =
\hbar^2k^2/2M = \pi^2\hbar^2/8MR^2$. (c) Con $\mu c^2 =
\qty{469}{MeV}$ y $R = \qty{2.1}{fm}$: $V_{0,\min} = \qty{23}{MeV}$.
(d) La pared $u(0) = 0$ hace del estado $s$ tridimensional el análogo de un
estado \emph{impar} unidimensional, que debe alojar un cuarto de onda dentro
del pozo: un pozo somero no puede, mientras que en una dimensión el estado
par sin nodos siempre liga.
\end{solution}

\begin{solution}{exo:b3:schrodinger-three-dimensions:10}
(a) $\kappa = \sqrt{2\mu c^2B}/\hbar c = \sqrt{2 \times 469 \times
2.22}/197.3 = \qty{0.231}{fm^{-1}}$: longitud de cola $1/\kappa =
\qty{4.3}{fm}$, el doble del alcance de la fuerza. (b) $k =
\sqrt{2\mu(V_0 - B)}/\hbar$ y $\cot kR = -\kappa/k$. (c) Iterando:
$kR = 1.83$, $k = \qty{0.87}{fm^{-1}}$, $V_0 - B = (\hbar ck)^2/\mu
c^2\cdot\tfrac12 \approx \qty{31}{MeV}$: $V_0 \approx \qty{34}{MeV}$.
(d) Dentro: $\int_0^R\sin^2kr\,\dd r \approx \qty{1.19}{fm}$;
fuera: $\sin^2(kR)/2\kappa \approx \qty{2.0}{fm}$: alrededor del $63\%$ de
la probabilidad está \emph{más allá} del pozo --- el deuterón habita sobre
todo la región en la que su fuerza de enlace ya se ha agotado, un halo
puramente cuántico.
\end{solution}

\begin{solution}{exo:b3:schrodinger-three-dimensions:11}
(a) $\dd\langle x\rangle/\dd t = \int x\,\partial_t\rho\,\dd^3r =
-\int x\operatorname{div}\vect\jmath\,\dd^3r = \int j_x\,\dd^3r =
\langle p_x\rangle/m$ (por partes, con los términos de borde anulándose).
(b) Derívese $\langle\vect p\rangle = -\iu\hbar\int\psi^*\vect\nabla
\psi$ y úsese la ecuación dos veces; los términos cinéticos se cancelan por
partes y queda $-\int|\psi|^2\vect\nabla V$. (c) Que $V$ sea a lo sumo
cuadrático: entonces $\langle\vect\nabla V(\vect r)\rangle = \vect\nabla V(\langle
\vect r\rangle)$ exactamente. (d) Detrás de una doble rendija, el paquete
son dos lóbulos; $\langle\vect r\rangle$ se desliza por el eje de simetría,
donde la partícula prácticamente nunca aterriza: el promedio describe el
centroide del conjunto, no la trayectoria de nadie.
\end{solution}

\begin{solution}{exo:b3:schrodinger-three-dimensions:12}
(a) Cinética: los gradientes de escala $1/a$ dan $\hbar^2/2m_{\text{e}}
a^2$; potencial: la nube se sitúa a distancias $\sim a$, lo que da
$-ke^2/a$ (los coeficientes exactos resultan ser $1$). (b) $\dd E/\dd
a = 0$ en $a = \hbar^2/m_{\text{e}}ke^2 = a_0$; $E(a_0) =
-ke^2/2a_0 = \qty{-13.6}{eV}$. (c) Por debajo de $a_0$, el coste cinético
crece como $1/a^2$, más deprisa que la ganancia $-1/a$: el principio de
incertidumbre grava la localización, y el átomo flota en el tamaño de
equilibrio. (d) Sustitúyase $ke^2$ por $Gm_{\text{e}}m_{\text{p}} =
\qty{1.0e-67}{J.m}$: $a_{\text{grav}} = \hbar^2/m_{\text{e}}G
m_{\text{e}}m_{\text{p}} \approx \qty{1.2e29}{m}$ --- mayor que
el universo observable: la gravedad es demasiado débil para cerrar una
órbita cuántica, y construye estrellas en vez de átomos.
\end{solution}

\begin{solution}{pb:b3:schrodinger-three-dimensions:1}
\textbf{1.} Onda $s$: $u = rR$ cumple la ecuación libre unidimensional
dentro de la esfera con $u(0) = u(R) = 0$: $u \propto \sin(\pi r/R)$, con
energía $\pi^2\hbar^2/2\mu R^2$.
\textbf{2.} La energía fundamental de una caja es positiva y crece al
encoger la caja --- localizar cuesta siempre energía cinética (principio de
incertidumbre); solo puede sumarse a la banda prohibida.
\textbf{3.} Con $\pi^2(\hbar c)^2/2\mu c^2 = \qty{3.76}{eV.nm^2}$:
$0.04$, $0.24$, $0.94$ y \qty{3.8}{eV} --- en $R \approx \qty{2}{nm}$
la corrección rivaliza con la propia banda prohibida.
\textbf{4.} La atracción baja la energía del par: la emisión se desplaza al
rojo. Va como $1/R$ frente al $1/R^2$ del confinamiento: para los puntos
pequeños gana el término de caja y la aproximación mejora.
\textbf{5.} $1240/1.74 = \qty{713}{nm}$: en el borde rojo de la visión.
\textbf{6.} Un solo color procedente de $10^{15}$ emisores certifica que la
síntesis los hizo todos del \emph{mismo tamaño} --- la monodispersidad, la
hazaña química que hay detrás de la física.
\textbf{7.} $1.97$, $2.34$ y \qty{2.70}{eV}.
\textbf{8.} Confinamiento $0.23$, $0.60$ y \qty{0.96}{eV}: $R =
\sqrt{3.76/\Delta E}$ da $4.0$, $2.5$ y \qty{2.0}{nm}.
\textbf{9.} Diámetros $8.1$, $5.0$ y \qty{4.0}{nm}: aproximadamente $13$,
$8$ y $7$ espaciados reticulares de anchura.
\textbf{10.} El punto azul es el más pequeño, y allí $E$ depende de $R$ con
la mayor pendiente: un error de un solo plano reticular desplaza más el
color --- los puntos pequeños son los menos indulgentes.
\textbf{11.} Derivando: $\dd E/\dd R = -2\Delta E_{\text{conf}}/R
= -\pi^2\hbar^2/\mu R^3$; en $R = \qty{2.5}{nm}$: $2 \times 0.60/2.5
= \qty{0.48}{eV/nm} = \qty{480}{meV/nm}$.
\textbf{12.} $\pm 5\% = \pm\qty{0.125}{nm}$: $\Delta E = \pm
\qty{60}{meV}$, es decir, $\Delta\lambda = \lambda^2\Delta E/hc \approx
\pm\qty{14}{nm}$ --- una dispersión de $\sim\qty{28}{nm}$, justo en la
tolerancia de la pantalla: un cinco por ciento es la \emph{frontera} de la
química aceptable.
\textbf{13.} Porque la pureza de color exige controlar el tamaño al nivel
de un solo plano reticular a lo largo de $10^{15}$ partículas --- el
crecimiento controlado que lograron Ekimov, Brus y Bawendi, y que citó el
comité Nobel de 2023.
\textbf{14.} Un punto solo puede emitir en su propia banda prohibida (la
más baja); absorber exige al menos esa energía, luego el bombeo debe ser más
azul. El exceso se relaja en vibraciones de la red: calor.
\textbf{15.} $1 - 460/630 = 27\%$ de la energía de cada fotón de bombeo
calienta la pantalla.
\textbf{16.} Por encima del nivel emisor, el espectro del punto es denso
(muchos estados de caja): la absorción tiene éxito en una banda ancha; la
excitación cae después al estado excitado más bajo, y la emisión ocurre
desde ese único nivel: estrecha.
\textbf{17.} Todos los puntos absorben con gusto el mismo ultravioleta
(banda ancha), mientras que cada tamaño emite su propio color: una lámpara,
muchas marcas --- lo inverso de la química de colorantes.
\textbf{18.} $\tfrac43\pi(4.0)^3/(0.3)^3 \approx \num{e4}$ átomos:
diez mil átomos que comparten un mismo conjunto de niveles discretos, tipo
hidrógeno --- ``átomo artificial'' está bien ganado.
\textbf{19.} La frontera del cristal: la pared confinante del grano
semiconductor sustituye a la atracción del núcleo como agente que retiene y
cuantiza al electrón.
\textbf{20.} $\pi^2(\hbar c)^2/2Mc^2R^2 = 384210/(2 \times 939
\times 9) \approx \qty{23}{MeV}$: la física nuclear habla en MeV porque las
cajas de femtómetros lo hacen.
\textbf{21.} Para un electrón, la estimación de caja se vuelve relativista;
honradamente, $E \sim \pi\hbar c/R \approx \qty{200}{MeV}$ --- y, sin
embargo, los electrones beta salen con unos pocos MeV: los electrones no
pueden ser constituyentes de los núcleos, el argumento que enterró el viejo
modelo nuclear con electrones y preparó la invención del neutrino.
\textbf{22.} $\pi^2\hbar^2/2ML^2 \approx \qty{2e-70}{J}$:
treinta y nueve órdenes por debajo del ruido térmico --- las personas no
están confinadas cuánticamente.
\textbf{23.} $\pi^2\hbar^2/2m_{\text{e}}R^2 = m_{\text{e}}c^2$ en
$R = \pi\hbar/\sqrt2\,m_{\text{e}}c \approx \qty{0.9}{pm}$, la
escala Compton: allí la energía de confinamiento basta para crear pares
electrón--positrón y la ecuación de una sola partícula abdica en favor de la
teoría cuántica de campos.
\textbf{24.} $E_{\text{conf}} \propto 1/MR^2$: las partículas más ligeras y
las cajas más pequeñas son más cuánticas, en una sola fórmula.
\textbf{25.} $E_{\text{g}} + \pi^2\hbar^2/2\mu R^2$ lleva $4.0 \to$
rojo, $2.5 \to$ verde y $2.0\,\unit{nm} \to$ azul: la partícula en una
caja, afinada por los químicos al nanómetro, ilumina el escaparate ---
cuantización por confinamiento convertida en electrónica de consumo.
\end{solution}
